\documentclass[11pt]{article}
\usepackage[T1]{fontenc}
\usepackage{lmodern}
\usepackage[margin=1in,top=.82in,bottom=.88in]{geometry}
\usepackage{amsmath,amssymb,amsthm,booktabs,microtype}
\usepackage[hidelinks]{hyperref}
\newtheorem{theorem}{Theorem}
\newtheorem{lemma}[theorem]{Lemma}
\newcommand{\ph}{\varphi}
\newcommand{\fl}[1]{\lfloor #1\rfloor}
\newcommand{\fr}[1]{\{#1\}}
\title{Three gap sequences from the Wythoff and Stolarsky arrays}
\author{Blake Stone mathematics repository\\[-2pt]\small Proof developed with OpenAI Codex}
\date{}
\begin{document}
\maketitle
\begin{abstract}
The second columns of the Wythoff and Stolarsky arrays differ by $0$, $1$ or $-1$.
Kimberling conjectured that the row indices for each value have exactly three
consecutive gap sizes. We determine the gaps and their limiting frequencies
from one first-return map of a golden rotation. The lower Wythoff map also
pairs the positive and negative lists in increasing order. Their gap
sequences agree term by term after relabeling.
\end{abstract}

\section{The three lists}
Let $\ph=(1+\sqrt5)/2$, $\beta=\ph-1$, $F_0=0$, $F_1=1$ and
$F_{i+2}=F_{i+1}+F_i$. We number both rows and columns from $1$.
The Wythoff array and the two Stolarsky entries needed here are
\begin{align*}
 W(n,k)&=F_{k+1}\fl{n\ph}+(n-1)F_k,\\
 S(n,1)&=\fl{(n-1)(\ph+1)+1+\ph/2},&
 S(n,2)&=\fl{\ph S(n,1)+1/2}.
\end{align*}
Set $D(n)=W(n,2)-S(n,2)$ and $R_s=\{n\ge1:D(n)=s\}$.
For example, $R_0$ starts $1,4,6,7,9$,
$R_+$ starts $2,5,10,13,15$ and $R_-$ starts $3,8,16,21,24$.
These give three gap questions in the OEIS \cite{oeis}.
The entry A387349 also points to Shallit's automata method \cite{shallit}.
We give a direct proof for the three sets, including occurrence, frequency and
the correspondence between two of the infinite gap sequences.

\begin{theorem}\label{main}
The possible gaps in each increasing list and the natural density of its row indices are
\[
\begin{array}{lccc}
\toprule
\text{Row indices}&\text{OEIS}&\text{successive gap sizes}&\text{row density}\\
\midrule
R_0 & \mathrm{A387349} & 1,2,3 & 1/2\\
R_+ & \mathrm{A387350} & 2,3,5 & \beta/2\\
R_- & \mathrm{A387351} & 3,5,8 & \beta^2/2\\
\bottomrule
\end{array}
\]
Within each gap sequence, the proportions for the short, middle and long gaps
are $\sqrt5-2$, $5-2\sqrt5$ and $\sqrt5-2$, respectively.
If $m$ is the $i$th member of $R_+$, then $\fl{m\ph}$ is the $i$th
member of $R_-$. This changes the successive gap lengths
$2,3,5$ into $3,5,8$, respectively.
\end{theorem}

\section{Recovering each sign from a single phase}
Write $f_n=\fr{n\beta}=\fr{n\ph}$, $a=\fl{n\ph}$,
$p=\beta^2=1-\beta$, $q=\beta^3=2\beta-1$, and $L=1/2$.
The identities
\[
 S(n,1)=n+a-\mathbf1_{f_n<\ph/2},\quad W(n,2)=2a+n-1,\quad
 (n+a)\ph=2a+n+p f_n
\]
follow by substituting $n\ph=a+f_n$ and $\ph^2=\ph+1$.
If $f_n<\ph/2$, the final rounding adds
$\fl{p f_n+1/2-\ph}$ to $2a+n$. This floor changes from $-2$ to $-1$
at $f_n=\beta/2$, since $p\beta/2+1/2-\ph=-1$.
If $f_n>\ph/2$, the final rounding adds $\fl{p f_n+1/2}=0$.
The phases for $n\ge1$ avoid $0,\beta/2,\ph/2$ by irrationality of
$\beta$ and $(2n-1)\beta$. Thus
\[
\begin{array}{c|ccc}
\text{row set}& R_+ & R_0 & R_-\\ \hline
 f_n & (0,\beta/2) & (\beta/2,\ph/2)&(\ph/2,1).
\end{array}
\]
Equivalently $D(n)=2\fl{n\beta}-\fl{(2n-1)\beta}$, as one sees
by taking the floor of $2f_n-\beta$.

\section{First return to a Fibonacci-scaled interval}
View real numbers modulo $1$. For $j\ge0$ and $c\in\mathbb R$ take the interval
parametrized by $c+(-\beta)^j x$ for $0\le x<L$. A visit at time $n$ has
$n\beta=h+c+(-\beta)^j x$ for some integer $h$.

\begin{lemma}\label{induced}
Let $u=F_{j+2}$ and $v=F_{j+3}$. Starting at a visit with coordinate $x$,
the first positive return time and its new coordinate are
\[
\begin{array}{lll}
\toprule
\text{starting interval}&\text{time}&\text{new coordinate}\\
\midrule
p\le x<L        &u&x-p\\
0\le x<L-q      &v&x+q\\
L-q\le x<p      &u+v&x-p+q\\
\bottomrule
\end{array}
\]
\end{lemma}
\begin{proof}
The identity $F_{t+1}\beta-F_t=(-\beta)^t\beta$ follows by induction.
The steps $(u,F_{j+1})$ and $(v,F_{j+2})$ therefore move the normalized
coordinate by $-p$ and $q$. Their determinant has absolute value $1$ by
Cassini's identity. Any return at time $d$ with starting and ending
coordinates $x,y$ hence has integers $a,b$ satisfying
\[
 d=au+bv,\qquad y-x=-ap+bq,\qquad |y-x|<L.
\]
If $0<d<u+v$, both coefficients cannot be positive, nor can both be
nonpositive. Opposite nonzero signs force $|y-x|\ge p+q>L$.
The remaining possibilities are positive multiples of just one vector.
A multiple $a\ge2$ gives $ap\ge2p>L$ and a multiple $b\ge2$ gives
$d\ge2v>u+v$. Thus the only possible early returns are $u$ and $v$.
The former stays in the interval exactly when $x\ge p$, and the latter
exactly when $x<L-q$. These conditions are disjoint. In the remaining
interval the combined step gives $0\le x-p+q<L$, since $0<p,q<L$.
\end{proof}

For the three row sets the parameters are
\[
\begin{array}{c|ccc}
 &j&c&x\text{ at phase }f\\ \hline
R_0&0&\beta/2&f-\beta/2\\
R_+&1&\beta/2&(\beta/2-f)/\beta\\
R_-&2&\ph/2&(f-\ph/2)/p.
\end{array}
\]
Neither endpoint of any of the three oriented intervals occurs at a positive
time. The interior cuts $x=p$ and $x=L-q$ also never occur, since adding
$u$ or $v$ would take them to an endpoint. The lemma gives the successive
gap sizes in Theorem~\ref{main}. For explicit witnesses, the starting rows
for the short, middle and long gaps are respectively $(6,4,1)$ in $R_0$,
$(13,2,5)$ in $R_+$ and $(21,3,8)$ in $R_-$.

\section{Frequencies and the pairing}
Irrational rotation equidistributes the phases $f_n$ in $[0,1)$.
Consequently a normalized subinterval of length $t$ in the chart at level
$j$ is visited with density $\beta^j t$. The whole chart has length
$\beta^j L$, giving the three row densities. Division by this density makes
the relative gap frequencies equal to
\[
 (L-p)/L,\quad (L-q)/L,\quad (p-L+q)/L.
\]
All three are positive. Substituting $p,q,L$ gives the announced constants.

For the bijection put $n=\fl{m\ph}$ with $m>0$. Multiplying $m\ph=n+f_m$
by $\beta$ gives $f_n=1-\beta f_m$. Hence $m\in R_+$ exactly when
$n\in R_-$, since $1-\beta^2/2=\ph/2$. Conversely, for $n\in R_-$, set
$m=\fl{n\beta}+1$. Then $m\ph=n+\ph(1-f_n)$ and
$0<\ph(1-f_n)<\ph(1-\ph/2)=\beta/2<1$. Thus $\fl{m\ph}=n$, and
the map is strictly increasing. The two normalized coordinates agree,
\[
 \frac{f_n-\ph/2}{p}=\frac{\beta/2-f_m}{\beta}.
\]
Their next returns use the same line of Lemma~\ref{induced}, which proves the
termwise correspondence and completes the theorem.

\paragraph{Verification.} The accompanying Lean project derives the three
sets from the displayed array definitions. It proves the first-return
lemma at every level, the exact gap sets, the frequencies and the increasing
pairing. The integer equidistribution argument is included. It first
telescopes counts for endpoints $1-\{k\beta\}$, then uses their density
to obtain arbitrary intervals.

\begin{thebibliography}{9}
\bibitem{oeis} C. Kimberling, Wythoff--Stolarsky difference and positions of its
three values, OEIS entries \href{https://oeis.org/A387348}{A387348},
\href{https://oeis.org/A387349}{A387349}, \href{https://oeis.org/A387350}{A387350}
and \href{https://oeis.org/A387351}{A387351}, 2025, with subsequent comments.
\bibitem{shallit} J. Shallit, \emph{An {“Experimental Mathematics”} Approach to Stolarsky Interspersions
via Automata Theory}, arXiv:2502.03312v2, 2026.
\href{https://arxiv.org/abs/2502.03312}{arXiv record}.
\end{thebibliography}
\end{document}
